\section{Death to the World}

\subsubsection{In Memoriam}

In the USA, today is Memorial Day, when the war dead are memorialized. By extension, it is also a time to remember the dead among our family and friends; it is the secular equivalent to All Soul's Day for an irreligious nation. We can also use the occasion to meditate on our own deaths, something we forget to do. Forgetting, sleeping, and death are all manifestations of a similar process.

With my own struggles with illness fresh in my mind, the death yesterday of Dennis Hopper is particularly memorable. For this writer, Mr. Hopper was at one point in his life what is called a ``role model''. I want to use E. F. Schumacher's\footnote{\url{http://www.smallisbeautiful.org}} ``Four Fields of Knowledge'' to explore what that could mean. First of all, there is the I and the not-I, or ``world'', or even ``you''. Those are both the subjects and objects of knowledge. Thus our knowledge derives from two sources: ``Inner Experience'' and ``Outer Appearance''. This leads to the schematization in Table \ref{tab:SchumacherFields}.

\begin{table}[h]\centering\small
\begin{tabular}{rcc}\toprule
& \textbf{Source} & \\\midrule
I & I --- inner & I --- outer \\
You/It & You/It --- inner & You/It --- outer \\\bottomrule
\end{tabular}
\label{tab:SchumacherFields}
\end{table}

This leads us to four questions:

\begin{enumerate}


\item Who am I?

\item What is the inner experience of others?

\item How do I appear to others?

\item What is happening in the external word?
\end{enumerate}


Question 2 relates to empathy and 4 leads to science, but for our purposes, I will explore 1 and 3.

Dennis Hopper succeeded at level 3: he was known --- at least in his exteriority --- to millions. In his breakout move, \emph{Easy Rider}, he and Peter Fonda succeeded in creating an image that many in that generation aspired to copy. Leaving aside the wisdom of roaming aimlessly on the highways on a motorcycle, or the use of mind-altering chemicals, it was still an image that captured the imaginations of millions, even if very few could succeed in pulling it off. But it was all exteriority, a pose, and it seems that Mr. Hopper remained in that mode his entire life. It was a life tied to the sensuous world and only the few with the looks, charm, and money can master that world, at least for a season, until the inevitable corruption arrives.

Some few men will turn aside from the sensuous world and seek to know themselves, not as another object in the world, but as the subject and agent of their own life. Then question 3 becomes irrelevant (except for those who play the guru game), and the focus is on question 1. Even for most of them, it is a transitory, not a permanent, condition. Therefore, it is important to remember to turn aside from the sensuous world to the higher world.

Yesterday I took a walk along the beach and did the following exercise. I recalled that state of consciousness called dispassion, or \emph{apatheia}, and then as I observed the things in the world, I would begin to detach from them. There may be a short twinge of sadness or regret at having to leave these things behind, but that passes as the sense of detachment grows. This is called dying before you die.


\flrightit{Posted on 2010-05-31 by Cologero}

